\documentclass{handout}
\usepackage{handoutSetup}\usepackage{handoutShortcuts}
\usepackage{subfigure}


\begin{document}
\handoutHeader


\begin{verbatimwrite}{\jobname.title}
CRRA Portfolio Choice with Two Risky Assets
\end{verbatimwrite}
\medskip\medskip\medskip

\handoutNameMake
\begin{verbatimwrite}{./body.tex}

\cite{merton:restat} and \cite{samuelson:portfolio} study 
  optimal portfolio allocation for a consumer with Constant Relative Risk
  Aversion utility $\uFunc(c) = (1-\CRRA)^{-1}c^{1-\CRRA}$ who can choose
among many risky investment options.

Using their framework, here we study a consumer who has wealth
$\aRat_{t}$ at the end of period $t$, and is deciding how much to
invest in two risky assets with lognormally distributed return factors
$\Risky_{t+1}=(\Risky_{1,t+1}, \Risky_{2,t+1})'$, $\log \Risky_{t+1} =
\risky_{t+1}=(\risky_{1,t+1}, \risky_{2,t+1})' \sim \left (
  \mathcal{N}(\risky_1,\sigma^{2}_{1}), (
  \mathcal{N}(\risky_2,\sigma^{2}_{2}) \right)'$, with covariance
matrix
\begin{equation*}
 \left(\begin{array}{ccc}\sigma_1^2 & \sigma_{12} \\ \sigma_{12}& \sigma_2^2\end{array}\right).
\end{equation*}


If the period-$t$ consumer invests proportion $\riskyshare_i$ of $\aRat_{t}$ in  risky asset $i$, $i=1,2$ (so that $\riskyshare_{1}=(1-\riskyshare_{2})$ and vice-versa), spending all available resources in the last period of life\footnote{The portfolio allocation solution obtained below induces back to earlier periods of life, as \cite{samuelsonFallacy,samuelsonJudgment} famously emphasized.} $t+1$ will yield:
\begin{equation}\begin{gathered}\begin{aligned}
c_{t+1}  =  \underbrace{(\riskyshare \cdot \Risky_{t+1} )}_{\equiv \Rport_{t+1}} \aRat_{t} \notag \label{eq:RportDef}
\end{aligned}\end{gathered}\end{equation}
where $\Rport_{t+1}$ is the portfolio-weighted return factor.  

\cite{cvAppendix} point out that a good approximation to the portfolio rate of return is obtained by
\begin{equation}\begin{gathered}\begin{aligned}
\rport_{t+1} = \risky_{1,t+1}+\riskyshare_{2} (\risky_{2,t+1}-\risky_{1,t+1})+ \riskyshare_2 (1-\riskyshare_2) \eta/2
\end{aligned}\end{gathered}\end{equation}
where
$$\eta=(\sigma_1^2+\sigma_2^2-2\sigma_{12}). $$

Using this approximation, the expectation as of date $t$ of utility at date $t+1$
is:
\begin{equation}\begin{gathered}\begin{aligned}
  \Ex_{t}[\uFunc(c_{t+1})] & \approx  (1-\CRRA)^{-1}\Ex_{t}\left[\left(\aRat_{t}e^{\risky_{1,t+1}}e^{\riskyshare_2 (\risky_{2,t+1}-\risky_{1,t+1})+\riskyshare_2(1-\riskyshare_2)\eta /2}\right)^{1-\CRRA}\right] \notag
\\                      & \approx  \underbrace{ (1-\CRRA)^{-1}\aRat_{t}^{1-\CRRA}}_{\text{constant $< 0$}}\underbrace{e^{ (1-\CRRA)\riskyshare_2(1-\riskyshare_2)\eta/2}\Ex_{t}\left[e^{(\risky_{1,t+1}+\riskyshare_2 (\risky_{2,t+1}-\risky_{1,t+1}))  (1-\CRRA)}\right]}_{\text{excess return utility factor}}
 \label{eq:exputil}
  \end{aligned}\end{gathered}\end{equation}
where the first term is a negative constant under the usual assumption that relative risk aversion $\CRRA>1.$

  Our foregoing assumptions imply that $$(1-\CRRA) (\riskyshare_1 \risky_{1,t+1}+\riskyshare_2 \risky_{2,t+1})
   \sim \mathcal{N}( (1-\CRRA)(\riskyshare_1\risky_1+\riskyshare_2 \risky_2),      (1-\CRRA)^2 (   \riskyshare_1^2\sigma_1^2+\riskyshare_2^2\sigma_2^2+2\riskyshare_1\riskyshare_2\sigma_{12}))$$
(using \LogELogNormTimes).  With a couple of extra
  lines of derivation we can show that the log of the expectation in \eqref{eq:exputil} is
\begin{equation}\begin{gathered}\begin{aligned}
  \log \Ex_{t}\left[e^{(\risky_{1,t+1}+\riskyshare_2 (\risky_{2,t+1}-\risky_{1,t+1}))  (1-\CRRA)}\right] & =  {(1-\CRRA)(\riskyshare_1 \risky_1+\riskyshare_2 \risky_2) + (1-\CRRA)^2 (   \riskyshare_1^2\sigma_1^2+\riskyshare_2^2\sigma_2^2+2\riskyshare_1\riskyshare_2\sigma_{12}))/2} \notag
\label{eq:Ex}
\end{aligned}\end{gathered}\end{equation}

Substituting from \eqref{eq:Ex} for the log of the expectation in \eqref{eq:exputil},
the log of the `excess return utility factor' in \eqref{eq:exputil} is
\begin{equation*}
  (1-\CRRA)\riskyshare_2(1-\riskyshare_2)\eta/2+(1-\CRRA) (\risky_1+\riskyshare_2(\risky_2-\risky_1))+(\CRRA-1)^2 (\sigma_1^2+\riskyshare_2^2 \eta+2\riskyshare_2(\sigma_{12}-\sigma_2^2))/2
.
\end{equation*}
The $\riskyshare$ that minimizes this log will also minimize the level; the FOC
for minimizing this expression is
\begin{equation}\begin{gathered}\begin{aligned}
         (1-2\riskyshare_2)\eta/2+ \risky_2-\risky_1+(1-\CRRA) (\riskyshare_2 \eta+(\sigma_{12}-\sigma_1^2)) & =  0 \\
         \quad(\risky_2-\risky_1+\frac{\eta}{2})+(1-\CRRA) (\sigma_{12}-\sigma_1^2) & = \CRRA \eta \riskyshare_2.
         \label{eq:riskyshareMS}
\end{aligned}\end{gathered}\end{equation}

So
\begin{equation}\begin{gathered}\begin{aligned}
\riskyshare_2 & =  \left(\frac{\risky_2-\risky_1+\eta/2+(1-\CRRA)(\sigma_{12}-\sigma_1^2)}{\CRRA\eta}\right)
\end{aligned}\end{gathered}\end{equation}
and note that if the first asset is riskfree so that $\sigma_{1}=\sigma_{12}=0$ then this reduces to 
\begin{equation}\begin{gathered}\begin{aligned}
\riskyshare_2 & =  \left(\frac{\risky_2-\risky_1+\sigma^{2}_{2}/2}{\CRRA\sigma^{2}_{2}}\right) \label{eq:riskyshare2}
\end{aligned}\end{gathered}\end{equation}
but the log of the expected return premium (in levels) on the risky over the safe asset in this case is $\EpremLog \equiv \log \Risky_{2}/\Risky_{1} = \risky_{2}-\risky_{1}+\sigma^{2}_{2}/2$
(recalling that we have assumed $\sigma_{12}=\sigma^{2}_{1}=0$), so \eqref{eq:riskyshare2} becomes 
\begin{equation}\begin{gathered}\begin{aligned}
\riskyshare_2 & =  \left(\frac{\EpremLog}{\CRRA\sigma^{2}_{2}}\right)
\end{aligned}\end{gathered}\end{equation}
which corresponds to the solution obtained for the case of a single risky asset in \handoutA{Portfolio-CRRA}.


\end{verbatimwrite}
\input ./body.tex

\input handoutBibMake

\end{document}
